The attacks generated in Section~\ref{sec:attack_generation} are used to train a channel-wise classifier. The classifier serves two distinct purposes. First, it identifies measurement channels that are likely to be compromised. Second, its calibrated channel probabilities are used to construct a sensor subset for resilient state estimation.

\subsection{Channel-wise localization}
Let $\mathcal{T}\subseteq\{1,\ldots,m\}$ denote the attacked sensor set and define the safe-sensor indicator
\begin{equation}\label{equ:support}
q_j=\begin{cases}
0,&j\in\mathcal{T},\\
1,&j\notin\mathcal{T}.
\end{cases}
\end{equation}
The localization model is a vector-valued map
$g(\mathbf{y};\boldsymbol{\theta})\in(0,1)^m$, whose $j$th component estimates the probability that sensor $j$ is safe. In the numerical study, $g$ is a four-layer fully connected network with three ReLU hidden layers and a sigmoid output layer. Its binary decision is
\begin{equation}\label{equ:classification}
\hat q_j=\mathbb{I}\!\left\{g_j(\mathbf{y};\boldsymbol{\theta})\ge\tfrac12\right\}.
\end{equation}
Given training pairs $\{(\mathbf{y}^{(i)},\mathbf{q}^{(i)})\}_{i=1}^N$, the parameters are obtained by minimizing the channel-wise binary cross-entropy
\begin{align}\label{equ:MLP_loss}
J(\boldsymbol{\theta})=-\frac{1}{N}\sum_{i=1}^{N}\sum_{j=1}^{m}
\big[&
q_j^{(i)}\log g_j(\mathbf{y}^{(i)};\boldsymbol{\theta})\nonumber\\
&+(1-q_j^{(i)})\log(1-g_j(\mathbf{y}^{(i)};\boldsymbol{\theta}))
\big].
\end{align}
The generative model enriches the training set with attacks near the effectiveness and detectability boundaries in \eqref{equ:def_success_FDIA}.

\subsection{Calibrated safe-set pruning}
Let $\hat{\mathcal{T}}^c=\{j:\hat q_j=1\}$ be the set declared safe by the classifier, and let $\mathcal{G}$ denote the sigma-field generated by the complete classifier output and the resulting ordering information. A held-out calibration set is used to estimate, for each candidate $j\in\hat{\mathcal{T}}^c$, the conditional safety probability
\begin{equation}\label{equ:calibrated_safety}
p_j=\Pr\{q_j=1\mid\mathcal{G}\},
\end{equation}
where calibration may be implemented by conditioning on score bins and other features included in $\mathcal G$. The following analysis assumes that these conditional probabilities are calibrated and that the channel safety events are conditionally independent given $\mathcal G$. These assumptions can be relaxed by replacing the product below with a joint lower confidence bound.

Order the declared-safe indices so that
$p_{j_1}\ge p_{j_2}\ge\cdots\ge p_{j_{|\hat{\mathcal{T}}^c|}}$.
For a prescribed reliability $\eta\in(0,1)$, define
\begin{equation}\label{equ:trust_number}
l_\eta=
\max\left\{
l\in\{0,\ldots,|\hat{\mathcal{T}}^c|\}:
\prod_{r=1}^{l}p_{j_r}\ge\eta
\right\},
\end{equation}
with the empty product equal to one, and select
\begin{equation}\label{equ:pruning_alg}
\hat{\mathcal{T}}_\eta^c=\{j_1,\ldots,j_{l_\eta}\}.
\end{equation}
This construction differs from a count-based confidence bound: it certifies the particular prefix supplied to the estimator.

\begin{proposition}[Safe-set precision]\label{Prop:safe_precision}
Under the calibration and conditional-independence assumptions,
\begin{equation}\label{equ:precision_after_pruning}
\Pr\{\hat{\mathcal{T}}_\eta^c\cap\mathcal{T}=\varnothing\mid\mathcal G\}
=\prod_{r=1}^{l_\eta}p_{j_r}\ge\eta.
\end{equation}
Equivalently, the safe-set positive predictive value
\begin{equation}\label{equ:PPV_def}
\operatorname{PPV}_{\rm safe}
=\frac{|\hat{\mathcal{T}}_\eta^c\cap\mathcal{T}^c|}
{|\hat{\mathcal{T}}_\eta^c|}
\end{equation}
equals one with probability at least $\eta$ whenever $l_\eta>0$.
\end{proposition}
\begin{proof}
The selected set is clean precisely when all $l_\eta$ selected channels are safe. Conditional independence given $\mathcal G$ and \eqref{equ:calibrated_safety} give the product in \eqref{equ:precision_after_pruning}; the inequality follows from \eqref{equ:trust_number}. Taking expectations with respect to $\mathcal G$ also gives the unconditional bound $\Pr\{\hat{\mathcal{T}}_\eta^c\cap\mathcal{T}=\varnothing\}\ge\eta$.
\end{proof}

\subsection{Sparsity--observability--precision relation}
The next result identifies the admissible reliability levels for which pruning preserves observability.

\begin{proposition}[Precision and observability]\label{Prop:precision_observability}
Suppose $\hat{\mathcal{T}}^c$ contains at least $k_0$ candidates. The pruned set is observable under A2 and is clean with probability at least $\eta$ if
\begin{equation}\label{equ:eta_observability}
\eta\le\eta_{\rm obs}
\triangleq\prod_{r=1}^{k_0}p_{j_r}.
\end{equation}
Moreover, $l_\eta$ is nonincreasing in $\eta$.
\end{proposition}
\begin{proof}
Condition \eqref{equ:eta_observability} makes $l=k_0$ feasible in \eqref{equ:trust_number}; hence $l_\eta\ge k_0$, and A2 implies observability. Cleanliness follows from Proposition~\ref{Prop:safe_precision}. Finally, increasing $\eta$ can only remove feasible values from the maximization in \eqref{equ:trust_number}.
\end{proof}

Attack sparsity imposes a separate physical limit. Let $k=|\mathcal T|$ denote the realized number of attacked sensors. Exactly $m-k$ sensors are genuinely safe; hence, the existence of a clean observable subset requires
\begin{equation}\label{equ:sparsity_budget}
m-k\ge k_0,\qquad\text{or equivalently}\qquad k\le m-k_0.
\end{equation}
False-safe declarations cannot relax \eqref{equ:sparsity_budget}; they only reduce the probability that the selected set is clean. Thus, \eqref{equ:sparsity_budget} quantifies the sparsity--observability limit, while \eqref{equ:eta_observability} quantifies the precision required to realize that limit using a learned localization model.

\begin{remark}
The chi-square detector and the localization network quantify different events. Equation~\eqref{equ:detection_probability} describes the probability of a global alarm, whereas \eqref{equ:precision_after_pruning} describes the probability that the particular sensor subset supplied to the estimator is uncontaminated. Both quantities enter the resilient-estimation analysis, but they should not be conflated.
\end{remark}

\begin{algorithm}[h!]\label{Alg:ADL_AAG_P}
\SetAlgoLined
\textbf{Offline:}\\
\begin{enumerate}
\item Construct $M_1$ and $M_2$ from \eqref{equ:M1}--\eqref{equ:M2}.
\item Generate labeled attacks by solving \eqref{equ:training_GM}.
\item Train the channel-wise classifier using \eqref{equ:MLP_loss}.
\item Calibrate the conditional safety probabilities in \eqref{equ:calibrated_safety}.
\end{enumerate}
\textbf{Online:}\\
\begin{enumerate}
\item Evaluate $g(\mathbf{y};\boldsymbol{\theta})$ and form $\hat{\mathcal{T}}^c$.
\item Sort the candidates by calibrated safety probability.
\item Compute $l_\eta$ from \eqref{equ:trust_number} and return
$\hat{\mathcal{T}}_\eta^c$ from \eqref{equ:pruning_alg}.
\end{enumerate}
\caption{Attack localization and calibrated safe-set pruning}
\end{algorithm}
